%------------------------------------------
% Título: Introdução
%------------------------------------------
\section{Introdução}
\label{sec:introducao}


\todo[inline]{\textbf{NOTA}: O texto aqui apresentado contém apenas orientações para a escrita do seu texto de qualificação e DEVE SER REMOVIDO.}

\textcolor{red}{A função da introdução é apresentar à banca o tema da pesquisa, delimitar o problema investigado e demonstrar a sua relevância científica, tecnológica e/ou social, sendo finalizada com a descrição de como o texto está organizado. No exame de qualificação, a introdução deve permitir que a banca avalie, logo de início, se o problema é pertinente, se os objetivos são exequíveis no prazo restante do mestrado e se a pesquisa possui aderência à área de concentração e à linha de pesquisa do PPCA. Por se tratar de um Mestrado Profissional, é essencial evidenciar a aplicabilidade do trabalho: o contexto real (organização, setor ou comunidade) em que o problema se manifesta e o produto técnico ou tecnológico esperado como resultado (por exemplo, software, aplicativo, processo, protocolo, base de dados, relatório técnico). Apesar de ser a primeira seção, normalmente é a última a ser revisada, pois só apresentamos aquilo que já escrevemos.}

\textcolor{red}{1) Comece com uma sentença descrevendo o tema da pesquisa e o que você pretende investigar. 2) Na sequência, apresente brevemente o estado atual do tema e passe para a motivação da pesquisa. Motivos são argumentos, isto é, justificativas que indicam por que a pesquisa é necessária. Argumente, preferencialmente, com base em dados e evidências da literatura ou do contexto profissional em que o problema ocorre, indicando a lacuna que o trabalho pretende preencher. O ideal é ter de dois a três argumentos, sustentados por citações, que mostrem a relevância do tema para a ciência e o impacto esperado para a sociedade ou para o setor produtivo.}

\textcolor{red}{3) Explicite, então, o problema de pesquisa. No PPCA, vamos mesclar duas estratégias: apresentar o objetivo geral (verbo no infinitivo + complemento), desdobrado em objetivos específicos, e as perguntas de pesquisa que, se respondidas, atendem ao objetivo geral. Indique também a(s) hipótese(s), quando houver, e as contribuições esperadas, separando contribuições científicas do produto técnico ou tecnológico a ser entregue. 4) Finalize a introdução com a descrição das seções que compõem o texto.}

\textcolor{red}{Atenção com relação à estrutura explicada. De modo simples, ela propõe que você responda quatro perguntas: 1) Qual é o tema da pesquisa e o que você irá estudar dentro dessa temática? 2) Quais os motivos que justificam a pesquisa e qual lacuna ela preenche? 3) Quais são os objetivos, as perguntas de pesquisa e as contribuições esperadas, incluindo o produto técnico ou tecnológico? 4) Como o texto de qualificação está organizado?}

\todo[inline]{\textbf{NOTA}: Qual o problema que você está tentando resolver através da pesquisa? Em que contexto profissional ele ocorre? Quais as restrições envolvidas?}

O problema considerado neste trabalho é expresso na seguinte questão:

\todo[inline]{A pergunta de pesquisa precisa ser respondida para que o objetivo geral seja alcançado e será o principal critério para avaliar a dissertação na defesa. Elabore-a de forma clara, delimitada e verificável. Você pode ter uma ou mais perguntas, mas lembre-se de que precisará respondê-las até o final do mestrado.}

O objetivo geral deste estudo é
\todo[inline]{Preencher com texto iniciando com um verbo no infinitivo seguido de complemento. Exemplos: “Desenvolver sistema/aplicativo/aplicação/ferramenta…”, “Propor método/processo/modelo…”, “Avaliar a eficácia de…”, “Investigar validade…”, “Testar validade do modelo…”, “Testar a relação de associação entre as variáveis A e B”, “Testar a relação de determinação entre as variáveis A e B” etc.}

Para alcançar o objetivo geral, foram definidos os seguintes objetivos específicos:
\todo[inline]{Liste de três a cinco objetivos específicos, também iniciados por verbos no infinitivo, que representem etapas concretas e mensuráveis da pesquisa. Indique quais já foram concluídos até a qualificação e quais serão realizados até a defesa.}

Visando contribuir com $\dots$, o contexto deste trabalho
está situado no uso de metodologias e técnicas de $\dots$
\todo[inline]{Descreva a situação ou o contexto geral referente ao assunto em questão, com informações atualizadas que deem consistência ao trabalho, e o contexto de aplicação profissional da pesquisa. Adicionalmente, apresente uma breve descrição da solução proposta e do produto técnico ou tecnológico esperado.}

\textbf{Organização do Trabalho.} As demais seções deste texto de qualificação estão organizadas da seguinte forma:
%
nos \textbf{Fundamentos Teóricos}, são apresentados os conceitos abordados neste trabalho, especificamente: \textcolor{red}{XXX, YYY e ZZZ};
%
nos \textbf{Trabalhos Relacionados}, são analisados os trabalhos relacionados à solução proposta, evidenciando a lacuna que esta pesquisa pretende preencher;
%
na \textbf{Solução Proposta}, são descritas as etapas do método da solução proposta para \textcolor{red}{XXXX};
%
na \textbf{Avaliação Experimental}, são apresentados o planejamento da avaliação da solução proposta e os resultados preliminares obtidos até o momento;
%
e, por fim, nas \textbf{Considerações Finais}, apresentam-se as considerações parciais sobre as atividades já desenvolvidas e o cronograma das atividades a serem realizadas até a defesa da dissertação.
